\section{Conclusion}

%This paper has extended decentralized stability conditions for power systems by introducing a transformation that preserves decentralization while addressing the lack of sectoriality at low frequencies in standard converters. By switching from the rectangular to the polar admittance frame, improved low-frequency phase behavior is achieved, and the combination of both frames provides a balanced solution across frequencies.

This work revisited decentralized small gain–phase conditions for power systems and identified their inherent limitations when applied to grid-forming converters. To overcome this, a generalized loop-shaping transformation was introduced, combining the rectangular and the polar admittance frames. Not only this addresses the lack of sectoriality at low frequencies in converters, but also renders tighter bounds for the phase behavior across the whole frequency range.

The method was demonstrated on a grid-forming converter connected to an infinite bus, where virtual admittance filtering further reduced conservativeness, and on the IEEE 14-bus system, where it successfully identified unstable devices. While the results are promising, challenges remain, particularly when studying systems that include both grid-forming and grid-following converters, as well as in guaranteeing stability and sectoriality of the grid impedance. Future work will focus on systematic selection of transformation matrices and on reducing conservativeness by embedding network topology information into the method.

%This paper has proposed an extended version of decentralized stability conditions for power systems. While preserving the decentralized nature of the conditions, we introduced a transformation to overcome limitations inherent in a direct application of the original method. Specifically, the lack of sectoriality at low frequencies in standard converters is addressed by changing the coordinate representation from the rectangular admittance frame to a polar frame. In the polar frame, favorable phase behavior is obtained at low frequencies, though this advantage diminishes at higher frequencies. To address this, the two frames are merged, as demonstrated in the paper.

%Using a grid-forming converter as an example, we showed how the proposed method can be applied to study the stability of a converter connected to an infinite bus. In this example, we additionally employed virtual admittance filtering to achieve tighter phase bounds and thereby reduce the conservativeness of the method. Finally, the approach was tested on the IEEE 14-bus system with an unstable configuration, where it successfully identified the problematic device.

%Despite these promising results, several open questions remain. For instance, while the proposed transformation performs well for GFM converters, we observed difficulties when applying it to grid-following converters. Furthermore, the transformation does not necessarily guarantee sectoriality of the grid impedance, an aspect that requires further investigation, along with the development of more systematic approaches to selecting transformation matrices that ensure such properties. Another important source of conservativeness stems from the fact that no geographical information about the network topology is considered; the method assumes that any converter may be connected to any bus. Incorporating this structural information into the transformation could potentially reduce conservativeness and enable more practical applications of the stability method. These directions will be the subject of future work.